% Chapter 9 — Collateral, Margin, and Lending — MINSKY
% Implements the ratified Phase 1 Design Ruling in full. Budget: 11 pp (10 + 1 ratified, G6).

\section{Collateral, Margin, and Lending}\label{ch:collateral}

This chapter discharges the amended clauses of constitution \S4 and \S8, and \S6
for the agreement units it introduces.

\subsection{One representation, keyed by declared regime}

One representation serves every collateral, margin, and lending arrangement the
ledger records. Every (wallet, unit) balance is the signed coordinate vector
(owned, lent, posted) of Chapter~3; the outright holding is the degenerate case
with all mass on owned; and any unit --- cash, a bond, a structured note, a
listed option --- may sit on the non-owned coordinates. Which coordinate an
inflow writes is decided by the legal regime declared once on the governing
agreement unit, never by the asset class of what moves. Title transfer writes
owned, plus a claim and an obligation that are themselves units, and touches no
marker coordinate; a security interest writes marker mass on the posted plane
and never touches owned; settled-to-market margin writes owned with no
obligation at all. The rest of this chapter is that key, applied.

The agreement is a unit --- constitution \S6: every right and obligation the
legal contract creates is a unit. Its declared terms carry the regime
(settlement, financing, or custody); the eligibility schedule and the per-line
valuation percentages; the thresholds, deadlines, and close-out mechanics of
its margining; the trapping conditions its distributions carry; and its
watches --- the valuation watch, the income watches --- declared at
registration like any other unit's. The agreement unit itself is non-valued:
the claims and obligations it creates carry the value, and pricing the
governing shell beside them would count the same economic reality twice
(\S4).

A claim or an obligation is a valued unit held positive by the obligee and
negative by the obligor: one unit, one fact, both sides, and conservation makes
disagreement between them unrepresentable. An obligation carries, as declared
terms, a deadline, a discharge predicate, and a fallback (\S14.1); by its
deadline it is discharged, compensated, or declared defaulted --- a duty that
fails to fire is a visible, overdue item, never silence --- and, like every
unit, it retires only from the zero vector.

Eligibility and haircuts are declared data, checked at the door --- the
Transaction Executor's admission check --- and never typed. Collateral is a
role a balance plays under an agreement, not a kind of thing: the type system
knows planes and agreement references; what may be posted, and at what
valuation percentage, is a schedule in the agreement's terms. Admitting a new
eligible asset class is a terms version, not a code change; an ineligible
posting is refused at the door. What the type system does enforce is stronger
than any schedule: a move writes exactly one coordinate of one unit, positive
quantity, the same coordinate debited at the source and credited at the
destination, and a move whose coordinate is not owned carries a mandatory
reference to its governing agreement. Orphan collateral mass is not refused;
it is unwritable.

\begin{lstlisting}
newtype AgreementRef = AgreementRef UnitId  -- the governing agreement, itself a unit
data Plane = Lent | Posted                  -- the marker planes

data Placement                              -- where a move writes
  = OnOwned                                 -- ownership moves; no agreement needed
  | OnMarker Plane AgreementRef             -- marker mass: unwritable without its agreement

data Regime = Settlement | Financing | Custody  -- declared once, on the agreement unit
\end{lstlisting}

Everything computable from the coordinates --- available quantity, encumbered
quantity, mass in possession, mass re-used --- is a projection, computed when
needed and never stored (\S4). Chapter~12 builds the reporting projections on
exactly this surface.

\subsection{The two guards}

Two guards govern the marker planes, and they are never collapsed into one
word. Both are \emph{used} throughout this chapter; their formal statements and
proof obligations are catalogued in Chapter~14.

\textbf{Coverage is an invariant.} No wallet delivers what it neither owns nor
holds. Checked at the door on every admission, per wallet and unit, net over
agreements:
\[
\sum_{G}\ \mathrm{posted}_G(w,u)\ \le\ \max\bigl(\mathrm{owned}(w,u),\,0\bigr),
\]
and the same bound with lent in place of posted. The sum runs over the signed
balances, so mass received offsets mass re-posted: a wallet that received 100
under one agreement and re-posted 60 under another nets to $-40$, within the
bound --- re-use of received collateral is thereby bounded by possession. The
$\max$ makes the bound bind only the posting direction: a short wallet with
nothing posted is admissible --- the issuer wallet that conservation itself
creates is one --- while a poster with owned $-5{,}000$ attempting to post
$5{,}000$ is refused at the door. Coverage is never lawfully false, because a
transaction is atomic and possession is instantaneous within it; that is what
makes it an invariant.

\textbf{Sufficiency is an obligation.} Collateral value at or above exposure
is lawfully false intraday: a market move, a knock, a call in flight all leave
it false with no party in breach. It is therefore an obligation, never an
invariant: a deadline, a discharge predicate (the delivery amount met), and
close-out as the declared compensation. A book below its sufficiency line is
not a broken state; it is an open, deadline-bearing item on the record, and
the deadline covers it.

\subsection{Title transfer: owned re-books; the claim is a unit}

Owned must be the coordinate the move machinery can act on. The taker of
title-transfer collateral may lawfully sell it outright; a design that keeps
the asset on the poster's owned makes that lawful sale a special constructor
and makes the manufactured payment --- the income the taker owes back to the
poster --- unrepresentable. So the poster's owned re-books to the taker, and
what the poster keeps is what the poster legally has: in the same transaction
the poster gains a \textbf{claim-for-equivalent} unit, and the taker's
negative balance in that unit is the redelivery (or repurchase) obligation.
Three rules complete the shape; each is a correctness requirement.

\emph{Pricing.} The claim is priced through the ordinary pricing layer,
parameterised --- as valuation under \S8 already is --- by state: here the
taker's. While the taker performs, the claim tracks the referenced asset (the
accrued redelivery value) and bears the taker's credit; on the taker's default
it decouples from the referenced asset and becomes a recovery claim on the
taker's estate. It is not ``priced identically to the asset'': that rule would
mark the poster long at par on the one day the mark exists for.

\emph{Identity.} Claims are per-agreement units --- per netting set. Close-out
nets within a master agreement and never across masters, and because owned
moves carry no agreement reference, the claim's netting-set identity cannot
live on a move: it lives in the unit's own declared terms, which name both the
referenced asset and the governing agreement. The unit registry grows by one
claim per financing transaction. Accepted: fungible claims would break
netting-set attribution at close-out, which is the one computation the claims
exist to get right --- and the counterparty-state pricing above is then
natural, because the claim is a per-counterparty object.

\emph{Timing.} Owned re-books at instruction --- trade-date booking --- and
the instruction-to-settlement gap is carried explicitly, never elided. Where a
record date falls inside the gap, the paying agent pays the record holder ---
still the deliverer --- while the ledger's owned shows the buyer. The
reconciling leg is a \textbf{market-claim obligation unit} created over the
gap: first-class, recorded, deadline-bearing --- an instrument, not a break.
Settlement state is thereby load-bearing for entitlement routing: a
correctness requirement of this design, not a reporting axis. Chapter~11
carries the settlement-side machinery; the mechanism is stated here because
the re-booking rule creates the gap it must close.

Cash under title transfer is the same shape, by \S4's own admission test: a
quantity earns a coordinate only when a distinct real-world action can change
that coordinate independently of ownership --- and taking cash under title
transfer \emph{is} the change of ownership, so it earns no marker coordinate.
The receiver owns the cash, reinvests it, bears the result, and owes an
equivalent-return debt: \S8's second case, the obligation itself a unit. A
custody marker over commingled cash would assert a per-dollar fact no bank
statement can confirm --- the nostro observes only the sum --- and its
intraday divergence from the statement is precisely the internal
reconciliation failure the ledger exists to abolish. The return obligation is
valued at the inflow amount --- par plus accrued at the declared rate --- so
net owned value cannot move at receipt: \S8's deposit-neutrality sentence
holds by construction, not by formula. Balance-sheet presentation, encumbrance,
and re-use figures are projections over claims, regimes, and coordinates ---
Chapter~12; no second store.

\subsection{Pledge: marker mass under the coverage bound}

A pledge changes possession independently of ownership, so it writes the
marker plane and only the marker plane: one move on the posted plane ---
poster $+Q$, taker $-Q$, the taker's negative balance projecting as received
as collateral --- admissible under the coverage invariant, with owned
untouched. That is the legal fact: the pledgor still owns. Because posting
never writes owned, $V=\sum \mathrm{owned}\cdot P$ is unchanged by any
post-and-return round trip, and \S4's ``only the owned coordinate carries
economic value'' becomes a machine-checkable property: insert a
post-and-return pair within any watch-free interval of the referenced
agreement and every later valuation is identical. The side condition is real
--- an insertion spanning the agreement's valuation watch lawfully changes
later firings and cash --- and the property, with it, is catalogued in
Chapter~14.

Re-use of pledge-form collateral is a posting of mass held on the received
ray, within the same coverage bound. Every re-pledge move carries a
source-agreement reference beside its governing reference --- the source is an
observable fact of the custody instruction, so nothing fictitious is stored
--- and re-use therefore projects as per-(unit, agreement) actuals in
Chapter~12.

\subsection{Cash margin: four cases, one machine}

Cash margin is four cases, keyed by the declared terms of the agreement unit,
and the enumeration is total: a margin flow whose agreement declares no regime
has no admissible booking and is refused at the door --- the system stops
rather than improvises (\S2).

\begin{center}
\begin{tabular}{@{}p{0.34\textwidth}p{0.17\textwidth}p{0.38\textwidth}@{}}
\toprule
The agreement declares & The inflow is & The transaction writes \\
\midrule
Settled-to-market (STM) variation margin & settlement (\S8 case 1) & owned; no obligation --- the day's exposure is extinguished \\
Collateralised-to-market (CTM) or title-transfer credit support & financing (\S8 case 2) & owned, plus a return-obligation unit accruing the declared rate \\
Security interest without right of use (segregated) & custody (\S8 case 3) & posted-plane marker mass; owned untouched \\
Security interest with right of use & financing upon exercise of the right & owned, plus the return-obligation unit, upon exercise \\
\bottomrule
\end{tabular}
\end{center}

The custody case raises no fungibility objection, because the segregated
account reconciles one-for-one against the marker. For cash under a right of
use, receipt into the receiver's own accounts is the exercise: commingled
right-of-use cash books as financing, never as custody --- the fourth row
collapses into the second at the moment commingling makes the custody claim
unverifiable.

\medskip\noindent\textbf{Episode (future): the settled-to-market margin day.}
FUT-IDX is a cash-settled listed future on index IDX, multiplier 10,
settled-to-market by the declared terms of the central counterparty (CCP)
agreement unit; wallet W-ALPHA is long one contract, and day~1 settled at
$1{,}000.00$.
\emph{Trigger.} The day-2 official settlement print, $990.00$, is recorded;
the Event Monitor emits the settlement event.
\emph{The contract.} The future's contract fires, reading the last applied
level, $1{,}000.00$, from PositionState (Chapter~6).
\emph{The transactions.} One transaction:
\begin{enumerate}[label=(\arabic*),nosep]
\item move owned(USD) $(1{,}000.00-990.00)\times 10 = $ \USD{100}, W-ALPHA
$\to$ W-CCP;
\item PositionState update: margin applied, settlement level $990.00$.
\end{enumerate}
No marker coordinate, no obligation unit: the payment settles the day's
exposure, and on unwind nothing is returned, which is correct. W-ALPHA's net
asset value falls by exactly 100, the day's loss.
\emph{The door.} Conservation per (unit, coordinate) on the paired-leg move;
an owned move, so no agreement reference is required; idempotence under the
cause-derived identifier --- a duplicate settlement event proposes nothing
new.
\emph{The contrast.} Had the CCP agreement declared CTM, the identical cash
move would be paired with a return-obligation unit valued at 100 and accruing
the declared rate: one declared term and one obligation unit of difference,
not a different machine.
\emph{The regime is a declared term of the agreement unit; the machine is the
same.}

\medskip
The regime is one declared bit carrying the largest balance-sheet consequence
in this chapter, and its misdeclaration is a class of error the recomputation
check of \S13 cannot catch: recomputation reproduces the declared bit and
faithfully reproduces the wrong answer. A margin agreement misdeclared as STM
omits the returnable-margin liability and its interest accrual; the converse
fabricates them. Detection is boundary reconciliation against the
counterparty's or the clearing house's margin statement --- nothing internal
detects it, and the reconciliation is therefore a named duty of the boundary,
not a redundancy. The repair is defined now, not improvised at the first
discrepancy: a ProductTerms amendment --- a new terms version declaring the
true regime --- \emph{plus} compensating transactions rebooking the obligation
units and accruals of the affected interval. Both, not one: the amendment
without the rebooking leaves the interval wrong; the rebooking without the
amendment repairs history and re-breaks tomorrow.

\subsection{Line valuation}

Collateral value under an agreement is the sum, over the lines the agreement
holds, of market value times the declared valuation percentage of that line.
The derivation is \S6 faithfulness: the ledger values what the agreement
declares, and the agreements declare lines --- an eligibility schedule with a
valuation percentage per line --- not floored packages. An agreement whose
declared terms define package-level valuation prices the package through the
ordinary pricing layer, admissible as declared data like any other term; and a
portfolio floor, where computed, is a projection --- never a stored valuation,
and never a haircut the agreement does not grant.

\subsection{Entitlements: determination and payment}

Entitlement \emph{determination} always reads the owned plane. The
instrument's contract is possession-blind and regime-blind: who is owed a
dividend, a coupon, a payout is a question about ownership, and ownership
lives on exactly one coordinate. Entitlement \emph{payment} routes through a
\textbf{conditional obligation unit} created by the determination, whose
discharge predicate carries the governing agreement's declared trapping terms
--- the sufficiency and default conditions under which the agreement withholds
a distribution. Where no agreement encumbers the owner's mass, the predicate
is vacuously true and the obligation discharges at once: in the ordinary state
the \emph{condition} vanishes, never the \emph{leg}. An unconditional
issuer-to-owner pass-through is forbidden: the agreements trap distributions
exactly when a delivery amount or a default stands, and \S6 binds the ledger
to the terms as declared.

\medskip\noindent\textbf{Episode (dividend): the unencumbered case.}
W-ALPHA holds $10{,}000$ ACME; the declared dividend is $1.50$ per share; the
record-date firing has already determined and recorded the entitlement,
\USD{15,000}, in PositionState (Chapters~5 and~6). Nothing encumbers
W-ALPHA's ACME.
\emph{Trigger.} The payment-date watch fires; the Event Monitor emits the
event.
\emph{The contract.} The stock's contract, on its third firing, pays against
the recorded entitlements.
\emph{The transactions.} One transaction:
\begin{enumerate}[label=(\arabic*),nosep]
\item the conditional payment obligation is created against the recorded
entitlement --- ACME's issuer wallet owes W-ALPHA \USD{15,000}; its predicate
carries the trapping terms of every agreement encumbering the holder's mass,
and there are none, so it holds vacuously;
\item move owned(USD) \USD{15,000}, ACME issuer wallet $\to$ W-ALPHA; the
obligation is discharged and retires at the zero vector.
\end{enumerate}
\emph{The door.} Conservation on both paired-leg moves; idempotence --- a
duplicate payment-date event finds the entitlement discharged and proposes
nothing; the obligation retires only from the zero vector.
\emph{In the ordinary state the condition vanishes, never the leg.}

\subsection{The knock while pledged}

The centrepiece assembles every rule above on one instrument. W-ALPHA owns one
one-touch option OT-1, written by W-BANK, paying \USD{1,000,000} on touch,
marked \USD{400,000} before the knock. It is pledged under security-interest
agreement G to counterparty wallet B at a declared valuation percentage of
50\%: W-ALPHA posted $+1$, B posted $-1$; collateral value to B
\USD{200,000} against an exposure of \USD{150,000}. The declared barrier is
$80.00$ on stock OMEGA; OMEGA prints an official close of $79.00$.

\begin{enumerate}[label=\textbf{\arabic*.},nosep,itemsep=2pt]
\item The low print is recorded; the Event Monitor emits the touch. The knock
is now on the record; no ledger state has changed --- the Monitor emits
events and cannot write state.
\item OT-1's contract fires on the event, and the Transaction Executor admits
its proposal: UnitStatus $\to$ \emph{triggered}; the entitlement is
\emph{determined} from the owned plane --- W-ALPHA; and the conditional
payment obligation is created --- W-BANK owes W-ALPHA \USD{1,000,000},
discharge predicate carrying G's declared trapping terms. Mass on the posted
plane is unchanged. A duplicate emission finds the unit triggered and proposes
nothing.
\item The pricing layer, state-parameterised, prices triggered OT-1 at 0.
Value has collapsed; mass is intact --- coordinates carry mass, prices carry
value.
\item G's contract, at its declared valuation watch, values B's received
projection: 0. The delivery amount is \USD{150,000}; it fires a margin-call
obligation with a deadline.
\item The payment obligation discharges under its predicate --- release only
to the extent no delivery amount stands. One transaction, two moves:
owned(USD) \USD{1,000,000} W-BANK $\to$ W-ALPHA, and posted(USD)
\USD{150,000} W-ALPHA $\to$ B under G, delivered as credit support and
discharging the call; \USD{850,000} stays free with W-ALPHA. Coverage holds
within the transaction: owned cash $1{,}000{,}000 \ge$ posted $150{,}000$.
Had G declared no trapping terms, the predicate would be vacuous and the full
payout would release; an agreement declaring proceeds attachment would have
G's contract post the cash instead --- declared terms decide.
\item The pledge closes: the marker return clears the posted plane of OT-1
(B/W-ALPHA $\pm 1 \to 0$); then the extinguishment move owned(OT-1) W-ALPHA
$\to$ W-BANK takes every wallet's vector for OT-1 to zero --- \textbf{the
ordering is forced by the coverage invariant}, which refuses extinguishing
owned while posted mass remains; and OT-1 retires at the zero vector.
\end{enumerate}

\emph{The door}, across the six steps: idempotence at step~2; conservation per
(unit, coordinate) on every move; coverage on step~5's transaction and again
as the ordering constraint of step~6; retirement only from the zero vector.
The residuals are named, not hidden. Between steps~2 and~5 the book is
lawfully under-collateralised: sufficiency is an obligation, and the deadline
covers it. The knock--revalue--call--discharge ordering is temporal, outside
the move algebra, and carries the safety and liveness model of Chapter~15.
And a worthless unit that no voluntary machinery will bury exits by the
supervised write-off path below.
\emph{One move-and-coordinate algebra, no special case: the agreement's
declared terms --- never the machinery --- decide where the payout lands.}

\subsection{The lent plane: securities lending}

The title-transfer derivation applies to the lending axis verbatim. Owned is
the coordinate the move machinery can act on, and the borrower's purpose is to
sell: a design that keeps the security on the lender's owned and marks the
loan as lent-plane mass makes the borrower's short sale either unrepresentable
--- the borrower has no owned balance to sell --- or coverage-violating ---
a lent-out marker standing against an owned balance the loan has emptied. A
securities loan that passes title therefore books as financing: owned re-books
at instruction; the lender holds the per-agreement claim-for-equivalent; the
borrower's negative balance in it is the redelivery obligation; income is
manufactured through the conditional obligation of the entitlement rule. The
lent plane itself remains what the posted plane is: the marker plane for
arrangements under which possession moves without ownership passing, written
under the same coverage bound. One derivation, two planes --- the regime key
applied to lending, not a second rule.

\medskip\noindent\textbf{Episode (securities loan): open, dividend, recall,
return.} Stock SIGMA trades at $40.00$; wallet W-LEND holds $1{,}000$ shares
and lends them to W-BORR under title-transfer lending agreement L, whose
declared terms carry the financing regime, an income watch, and the redelivery
window. The collateral leg of the loan follows the regime key of this chapter
unchanged and is not repeated here.

\emph{Trigger 1 --- open.} The executed loan instruction is recorded; L's
contract fires. One transaction:
\begin{enumerate}[label=(\arabic*),nosep]
\item move owned(SIGMA) $1{,}000$, W-LEND $\to$ W-BORR --- re-booked at
instruction;
\item issuance of the per-agreement claim-for-equivalent CLM-L: move
owned(CLM-L) $1{,}000$, W-BORR $\to$ W-LEND; the borrower's $-1{,}000$ is the
redelivery obligation, and CLM-L's declared terms name SIGMA and L.
\end{enumerate}
W-BORR may now sell outright --- the purpose of the borrowing. CLM-L is
priced by the borrower's state: it tracks SIGMA while W-BORR performs
(\USD{40,000} at inception) and decouples to a recovery claim on default.

\emph{Trigger 2 --- a dividend mid-loan.} SIGMA declares $2.00$ per share.
Determination reads the owned plane: W-BORR, holding $1{,}000$. On the payment
date, move owned(USD) \USD{2,000}, SIGMA issuer wallet $\to$ W-BORR. The same
recorded event fires L's contract on its income watch: the manufactured
payment is the conditional obligation --- W-BORR owes W-LEND \USD{2,000},
predicate carrying L's declared conditions (no default, no undischarged margin
deficiency). In the ordinary state it discharges at once: move owned(USD)
\USD{2,000}, W-BORR $\to$ W-LEND. The lender's cash offsets the dividend drop
in CLM-L's price: a stockholder's profit and loss, exactly, while the borrower
performs.

\emph{Trigger 3 --- recall.} W-LEND's recall notice is recorded at the
boundary; L's contract fires, and the redelivery obligation acquires its
deadline --- the declared redelivery window --- with redelivery of $1{,}000$
SIGMA as its discharge predicate and close-out compensation as its declared
fallback. By \S14.1 it is discharged, compensated, or declared defaulted by
the deadline; an unmet recall is a visible, overdue item, never silence.

\emph{Trigger 4 --- return.} The redelivery is recorded. One transaction:
\begin{enumerate}[label=(\arabic*),nosep]
\item move owned(SIGMA) $1{,}000$, W-BORR $\to$ W-LEND;
\item move owned(CLM-L) $1{,}000$, W-LEND $\to$ W-BORR, taking every wallet's
vector for CLM-L to zero; CLM-L retires at the zero vector, and the
redelivery obligation is discharged within its deadline.
\end{enumerate}

\emph{The door}, across the four triggers: conservation per (unit, coordinate)
on every paired-leg move; the claim retires only from the zero vector; the
deadline-bearing obligation is visible on the record until discharged.
\emph{A securities loan is title-transfer financing on the lending axis: the
same re-booking, the same claim, the same conditional leg --- no new
machinery.}

\medskip
The operational lifecycle of lending --- locates, buy-ins, and the
securities-financing reporting and settlement-discipline regimes that bind it
--- is the subject of the \emph{SBL Operations Companion}; the ledger shape is
complete here.

\subsection{Retirement, and the supervised write-off}

A lifecycle event extinguishes value, never mass (\S4). A knock, a default, an
expiry changes UnitStatus and price --- never a coordinate balance --- and a
unit leaves the ledger only from the zero vector, the agreement's machinery
clearing the marker planes first, in the order the coverage invariant forces.

Zero-vector retirement has no voluntary exit for a defaulted issuer's marker
mass: nobody returns a worthless unit, and dead units would otherwise
accumulate --- polluting the encumbrance projections and drowning the
obligation-liveness queue in items that can never discharge. The
\textbf{supervised write-off path} is the exit: an issuer-default lifecycle
event permits clearing the marker planes by compensating moves against
explicit, simultaneous loss recognition --- recorded, authorised, and open,
without the voluntary agreement machinery, and never silent. The centrepiece's
step~6 needs the voluntary machinery; a book that cannot run it runs this one.
